\documentclass[11pt,a4paper,onecolumn]{article}
\usepackage[T1]{fontenc}
\usepackage[utf8]{inputenc}
\usepackage{lmodern,microtype}
\usepackage[margin=1in]{geometry}
\usepackage{amsmath,amssymb,amsthm,mathtools}
\usepackage{booktabs,tabularx,array,enumitem}
\usepackage{float}
\usepackage{needspace}
\usepackage[section]{placeins}
\usepackage{tikz}
\usetikzlibrary{arrows.meta}
\usepackage[framemethod=default]{mdframed}
\usepackage{xcolor}
\usepackage[hidelinks]{hyperref}
\hypersetup{pdftitle={An L1-Authorized Sequencing Service with Regional Failover},
pdfsubject={Proposed protocol: model, algorithms, proofs and complexity}}
\usepackage[nameinlink,noabbrev]{cleveref}
\usepackage{fancyhdr}
\setlength{\headheight}{14pt}
\pagestyle{fancy}\fancyhf{}
\lhead{\small An authorized sequencing service}\rhead{\small Regional PBFT revision}
\cfoot{\thepage}
\setlength{\emergencystretch}{2em}
\setlist{nosep,leftmargin=*}
\newcolumntype{Y}{>{\raggedright\arraybackslash}X}
\newcommand{\Hash}{\mathsf H}
\newcommand{\Sig}{\mathsf{Sign}}
\newcommand{\Ver}{\mathsf{Verify}}
\newcommand{\QC}{\mathsf{QC}}
\newcommand{\RC}{\mathsf{RC}}
\newcommand{\Cfg}{\mathsf{Cfg}}
\newcommand{\gv}{v_{\mathrm G}}
\newcommand{\lv}{v_{\mathrm L}}
\newcommand{\PC}{\mathsf{PC}}
\newcommand{\AB}{\mathsf{AB}}
\newcommand{\NV}{\mathsf{NV}}
\newcommand{\botv}{\bot}
\newcommand{\pre}{\preceq}
\newcommand{\wt}{\operatorname{wt}}
\newcommand{\id}{\operatorname{id}}
\newcommand{\tip}{\operatorname{tip}}
\newcommand{\parh}{\operatorname{parent}}
\newcommand{\Tag}[1]{\mathsf{#1}}
\newcommand{\Proc}[1]{\textnormal{\textsc{#1}}}
\newtheorem{definition}{Definition}[section]
\newtheorem{assumption}[definition]{Assumption}
\newtheorem{lemma}[definition]{Lemma}
\newtheorem{theorem}[definition]{Theorem}
\newtheorem{proposition}[definition]{Proposition}
\theoremstyle{remark}\newtheorem{example}[definition]{Example}
\newtheorem{remark}[definition]{Remark}
\newcounter{protocol}
\newenvironment{protocol}[2]{\refstepcounter{protocol}\label{#1}\begin{mdframed}[nobreak=true,linewidth=.6pt,skipabove=10pt,skipbelow=10pt,innertopmargin=8pt,innerbottommargin=8pt]\small\raggedright\noindent\textbf{Protocol \theprotocol. #2}\par\smallskip}{\end{mdframed}}
\crefname{protocol}{Protocol}{Protocols}
\Crefname{protocol}{Protocol}{Protocols}
\crefname{definition}{Definition}{Definitions}
\crefname{assumption}{Assumption}{Assumptions}
\crefname{lemma}{Lemma}{Lemmas}
\crefname{theorem}{Theorem}{Theorems}
\crefname{proposition}{Proposition}{Propositions}
\title{An L1-Authorized Sequencing Service with Regional Failover\\
\large A PBFT Protocol for Verifiable Mini-Block Streams}
\author{}\date{28 September 2026}
\begin{document}
\begin{center}
{\Large\bfseries An L1-Authorized Sequencing Service\\with Regional Failover\par}
\vspace{.5em}
{\normalsize A PBFT Protocol for Verifiable Mini-Block Streams\par}
\vspace{.4em}{\small 28 September 2026}
\end{center}
\begin{abstract}
We specify an L1-authorized sequencing service whose producer rotates within
and between geographic regions. Regional PBFT commits yield certified previews;
explicit global agreement and FFG links yield durable delivery and geographic
activation. Producer push, peer pull and a cross-region backbone separate data
transport from confirmation. We give complete protocols, illustrative traces,
agreement and verification proofs, and latency and communication bounds for
normal service, local faults and regional outages. The regional normal path uses
quadratic message traffic without per-batch authenticated-broadcast completion.
View-change evidence and global report agreement are accounted for separately.
This is a proposed protocol;
consumer execution and publication lie outside its interface.
\end{abstract}

\section{Introduction}\label{sec:intro}
The purpose of the service is to deliver a continuous, verifiable transaction
stream under L1-authorized production rights. At any time one geographic region
is active, and one sequencer in that region produces ordered mini-blocks.
A local failure transfers production to another regional member; a regional
outage transfers it to another region. A proposer consumes the stream through
the same interface before and after these changes. 

\begin{figure}[H]\centering
\input{figures/service-architecture-pbft.tex}
\caption{The sequencing service and its external interfaces. The three drawn
sequencers illustrate communication roles; committee sizes and fault bounds
are specified in Section~\ref{sec:model}. A three-member equal-weight committee
under a strict one-third bound tolerates no Byzantine member.}
\label{fig:architecture}
\end{figure}

The active region confirms ordered mini-block batches; consumers receive leader
pushes and repair gaps by pulling from regional peers. Backbone synchronization
supplies the data and reports for global checkpoints. Local view changes preserve
the regional certified prefix. Geographic changes retain the global control
prefix, including every finalized delivery, and discard only its unretained
preview tail. A concrete synchronous report protocol precedes FFG voting,
making convergence and finality explicit.


\section{System Model and Setup}\label{sec:model}
\subsection{Participants, weights, and authorization}
The fixed sequencer set is $V=\biguplus_{i=1}^{m}V_i$, partitioned into $m$ regions.
Write $V_i=\{S_{i,1},\ldots,S_{i,n_i}\}$ and $N=|V|$. Sequencer $S_{i,j}$ has
weight $\omega_{i,j}>0$; equivalently $w_{S_{i,j}}=\omega_{i,j}$. For a set $U$,
\[
 \wt(U)=\sum_{S\in U}w_S,\qquad W_i=\wt(V_i),\qquad W=\wt(V).
\]
A quorum of committee $Q\subseteq V$ contains distinct members of weight
$>2\wt(Q)/3$. A strict threshold is used throughout. $P$ is an external
proposer consuming the stream; execution and publication by $P$ are outside
this ordering service. Clients submit transaction bytes to the active region.

L1 installs a configuration
\begin{equation}\label{eq:cfg}
 \Cfg_c=(c,V,\{V_i\},w,\mathcal L,\Theta,G_0).
\end{equation}
Here $c$ is the configuration identifier, $\mathcal L$ contains one ordered
cyclic list of every member of each region and one cyclic list of all regions,
and $G_0$ fixes the initial region, stream tip and finalized genesis checkpoint.
The parameter collection $\Theta$ contains regional and global delay bounds
$\Delta_R(i),\Delta_G$, regional retry duration $T_V(i)$, global epoch duration
$T_G$, regional batch limits $b_R,B_R$, global report limits $b_G,B_G$,
announcement period $T_A$, and complete request deadlines $\theta_R,\theta_G$.
The $b$ limits count mini-blocks or regional records as specified below; the
$B$ limits are bytes including the corresponding complete evidence.
Positive integers $K_G,N_G$ bound consecutive global no-progress epochs and
planned global tenure. All domains and size limits are fixed by $c$.
When the active region is clear, suppress the argument $i$ on $\Delta_R,T_V$.
L1 delegates local primary changes and finalized geographic changes through
this configuration; they require no additional L1 transaction.

\subsection{Cryptographic and communication interfaces}
$\Tag{Send}_{S\to T}(m)$ is an authenticated point-to-point transmission.
$\Sig_S(m)$ and $\Ver_S(m,\sigma)$ are signing and verification; honest
signatures are unforgeable. $\Hash$ is a collision-resistant binding hash of a
canonical, length-delimited encoding. Every signed tuple includes a type tag,
$c$, and its instance indices. Semantic identifiers exclude witness encodings
and signer subsets. No signature aggregation is assumed. $\botv$ denotes
absence of a value; it is not a mini-block.

The global layer uses synchronized epoch boundaries. Regional retries use
local elapsed-time timers at the same rate; a member never advances its timer
because of an unverified peer claim. Between responsive honest nodes, a
complete admitted object arrives and is verified within $\Delta_R$ regionally
or $\Delta_G$ globally. Actual regional delay is $\delta_R\leq\Delta_R$.
The bounds include retrieval of referenced bytes, queueing and verification,
including view-change bundles. They are workload/capacity assumptions,
not consequences of a short digest or a bounded number of messages.

\begin{assumption}[Faults and persistent state]\label[assumption]{ass:env}
The Byzantine set $B\subseteq V$ is fixed within $c$ and
$\wt(V_i\setminus B)>2W_i/3$ in every region. Historical honest keys remain
uncompromised. Honest members persist votes, retired views, prepared and commit
certificates, and authorizations before sending the corresponding messages.
Restart does not erase any signing restriction.
\end{assumption}
Honest nodes verify full bodies before PREPARE or global endorsement. They
retain and serve every prepared batch and its evidence while its sequence is
unresolved, and committed regional data and adopted global prefixes throughout
$c$. Unresolved evidence is kept across local timeouts; a geographic activation
may discard its unadopted data from the service output, but not signing history.

For progress we separately assume:
\begin{itemize}
\item A fixed responsive honest global set $H_G$ has $\wt(H_G)>2W/3$
throughout the analyzed execution and meets the global delay bound.
\item In a live regional interval a fixed honest set $H_i\subseteq V_i$ has
$\wt(H_i)>2W_i/3$ and meets the regional bound throughout entry and retries.
\end{itemize}
A healthy eligible successor $j$ has $\wt(H_G\cap V_j)>2W_j/3$ and the regional
bound on that set. This may be the recovered current region. Whole-region
outages are permitted only insofar as the global envelope still holds.
Local certificate safety does not require these timing bounds; end-to-end
progress additionally needs the capacity and tenure conditions stated later.

\subsection{Views, sequences, and heights}\label{sec:indices}
A global view $\gv$ authorizes one region $V_i$. Within it, $u\geq1$ is a
regional batch sequence number and $\lv\geq0$ is a PBFT retry view \emph{for
that sequence}. Thus the primary is indexed by $(\gv,u,\lv)$; $\lv$ resets to
zero when $u$ advances. Retries preserve $u$ and the committed predecessor.
For the configured local list $(L_{i,0},\ldots,L_{i,n_i-1})$, define
\begin{equation}\label{eq:primary}
 \ell_i(u,\lv)=L_{i,(u-1+\lv)\bmod n_i}.
\end{equation}
There is one outstanding sequence. The base primary rotates after each
committed batch, and each failed retry advances to the next member. This is
the local planned-rotation policy; there is no separate failed-GAP counter.
A new global view restarts $u=1,\lv=0$ at its inherited stream tip.

\begin{table}[htbp]\centering\small
\begin{tabularx}{\textwidth}{@{}p{.16\textwidth}Y@{}}\toprule
Symbol & Meaning\\\midrule
$i$ & Active geographic region; an identity, not a round.\\
$\gv$ & Global view, incremented at every finalized geographic activation,
including recovery in the same region.\\
$u$ & Batch sequence in $\gv$; advances only after a regional commit certificate.\\
$\lv$ & PBFT retry of $u$; increases on its timer while $u$ is unresolved.\\
$\alpha$ & Author view of a batch; remains fixed when a later primary re-proposes it.\\
$e$ & Global checkpoint epoch, occupying $[(e-1)T_G,eT_G)$.\\
$h$ & Mini-block height; continues from the inherited tip across activations.\\\bottomrule
\end{tabularx}
\caption{Indices for the regional PBFT and global control layers.}
\end{table}
Write $M_h$ for mini-block $h$, $q_h$ for its prefix digest, and $L\pre L'$
when stream $L$ is a prefix of $L'$. For example, retrying sequence $u=7$
from $\lv=0$ to $1$ changes its primary but not its parent or sequence number.
If a batch authored at $\alpha=0$ is retained, its author and bytes remain
unchanged. Committing it permits sequence $8$ to start in view zero.

\section{Service Specification}\label{sec:spec}
The service exports signed, ordered mini-blocks through
\[
 \Tag{Preview}_P(\gv,h,M_h,\pi_R),\qquad
 \Tag{Deliver}_P(h,M_h,\pi_G),
\]
where $\pi_R$ is regional commit evidence and $\pi_G$ is global finality evidence.
Raw producer pushes produce neither event. A preview is stable across local
PBFT retries and primary changes within its global view. A delivery is
append-only across geographic changes. A consumer persists its delivered log
$L_P$ and a separately labeled regional preview suffix. A geographic activation
may replace only the preview beyond the globally retained prefix.

\begin{definition}[Authorized, verifiable stream]\label{def:service}
Except with negligible cryptographic failure, the service requires:
\begin{enumerate}[label=\textnormal{(\roman*)}]
\item \emph{Validity and integrity.} Every output mini-block has its complete
body, matching commitment and authorized original producer. Each delivered
position is emitted once.
\item \emph{Agreement.} Honest consumers' delivered logs are pairwise
prefix-compatible at all times, and each grows monotonically. Regional
certified histories are prefix-compatible within one global view.
\item \emph{Progress.} In a live region, an outstanding sequence eventually
commits a valid batch. Under the global envelope, sufficient tenure and bounded
backlog, new mini-blocks are eventually adopted and delivered. After finitely
many unsuccessful geographic selections, a continuously healthy eligible
region resumes service, possibly through recovery in the same region.
\item \emph{Verification.} Complete correctly certified ranges pass the
specified predicates. Changed bytes, reordered or discontinuous ranges,
unauthorized producers and conflicting finalized histories are rejected.
\end{enumerate}
\end{definition}
Validity concerns ordered bytes and authority, not EVM execution. Each batch
contains at least one mini-block, but a mini-block may contain no transactions.
Progress therefore promises neither inclusion of a particular transaction nor
censorship resistance or ordering fairness. A previously prepared batch may
survive a primary change; an honest replacement is not entitled to overwrite it.

Regional decision latency is measured from a primary's complete proposal to
regional commitment; preview latency ends at the consumer's verified Preview.
Delivery latency ends at finalized Deliver. Recovery includes detection and
state transfer; planned handoff starts at its due boundary. Regional and
backbone communication, and body bytes and certificate bits, are counted
separately.

\section{Protocol Objects and Interfaces}\label{sec:objects}

\subsection{Mini-blocks and immutable batch values}
At the start of global view $\gv$, an activation proof $A$ fixes a tip $(h,q)$
and active region $i$. Genesis supplies the corresponding witness for view zero.
Define the initial regional record
\[
 R_0=((c,\gv,0),\Hash(A),\Tag{INIT},h,q).
\]
Its witness is the activation proof, not a regional vote. For sequence
$I=(c,\gv,u)$, let $p_u=\Hash(R_{u-1})$ be its certified predecessor identifier.
For an ordered transaction list $\mathcal X_h$, a batch authored in retry
$\alpha$ by $\ell_i(u,\alpha)$ contains
\begin{align}
 H_h&=(c,\gv,u,\alpha,h,q_{h-1},\Hash(\Tag{TXLIST},\mathcal X_h)),\label{eq:block}\\
 q_h&=\Hash(\Tag{MINIBLOCK},H_h),\qquad
 M_h=(H_h,\mathcal X_h,\Sig_{\ell_i(u,\alpha)}(\Tag{ORDER},H_h)).\nonumber
\end{align}
An empty transaction list has a canonical commitment. A proposal value is
\begin{equation}\label{eq:proposal}
 x=(I,\alpha,p_u,h_0,q_{h_0},h_1,q_{h_1}),\quad 1\leq h_1-h_0\leq b_R.
\end{equation}
Its witness contains every mini-block in the range and the original author's
signature $\Sig_{\ell_i(u,\alpha)}(\Tag{BATCH},I,\alpha,p_u,\Hash(x))$.
$\Proc{ValidBatch}(I,x)$ verifies this signature, the activated region,
certified predecessor, consecutive heights and parent digests, identical author
view throughout the batch, every ORDER signature and body commitment, and the
configured complete size bound $B_R$. The per-position checks are denoted
$\Proc{ValidMiniBlock}$. No execution or state-root check is implied.
A later primary may carry $x$ unchanged; its new signature belongs to the PBFT
wrapper and does not replace the original ORDER or BATCH signatures.

\subsection{PBFT certificates and view-change objects}\label{sec:pbftobjects}
For a statement $m$, $\QC_Q(m)$ contains signatures by distinct members of $Q$
of total weight $>2\wt(Q)/3$. $\Proc{VerifyQC}$ verifies each signature and this
threshold. For retry $v=\lv$, let $d=\Hash(x)$ and define phase statements
\[
 m_{\mathrm{phase}}=(\Tag{phase},I,v,p_u,d),\qquad
 \mathsf{phase}\in\{\mathsf{PREPREPARE},\mathsf{PREPARE},\mathsf{COMMIT}\}.
\]
A PREPREPARE is signed by $\ell_i(u,v)$. All members, including the primary,
send PREPARE under the same rule. A prepared certificate $\PC_v(x)$ consists
of that signed PREPREPARE and $\QC_{V_i}(m_{\mathrm{PREPARE}})$.
Its portable verifier checks the parent, indices, primary signature, quorum,
and batch witness. It need not contain the NEWVIEW proof previously checked
by honest prepare signers. The honest-supermajority assumption is essential
to this certificate abstraction.

The semantic regional outcome and its witness are
\begin{equation}\label{eq:rc}
 R_u=(I,p_u,d,h_1,q_{h_1}),\qquad
 \RC_u=(v,\QC_{V_i}(m_{\mathrm{COMMIT}}),x,\text{batch witness}).
\end{equation}
$\Proc{VerifyRC}$ checks the COMMIT quorum, $\alpha\leq v$, the complete valid
batch and its exact predecessor, and equality of the resulting $R_u$.
The committing retry $v$ and the signer subset are excluded from $R_u$'s
identifier. The same batch can commit in different retries without creating
different semantic outcomes. There is no regional GAP outcome: a timeout
retries the unresolved $u$ rather than certifying its absence.

For target retry $v\geq1$, a VIEWCHANGE from member $S$ is
\[
 \mathsf{VC}_{S,v}=(\Tag{VIEWCHANGE},I,v,p_u,\mathsf{pc}_S),
\]
with $S$'s signature. Here $\mathsf{pc}_S$ is its highest prepared certificate
from a retry $a<v$, including the full batch witness, or $\botv$ if none exists.
The certificate bytes are a detached witness to a signed semantic summary
$(a,\Hash(x))$; alternative signature subsets are the same report value.
A NEWVIEW witness $\NV_v$ is a set of such origin-signed reports of distinct
senders with weight $>2W_i/3$, all for $(I,v,p_u)$ and fully verified.
If any report is non-bottom, $\Proc{Select}(\NV_v)$ is the batch with the
highest prepared retry. Two different values at that same retry are invalid;
under the fault assumption they cannot both have prepared certificates.
If all reports are bottom, $\Proc{Select}(\NV_v)=\botv$ and the primary may
author a fresh valid batch at $\alpha=v$.
The signed NEWVIEW binds $(I,v,p_u,\Hash(\NV_v),\Hash(x))$ and carries a new
PREPREPARE for $x$. Receivers check the reported selection themselves.

\subsection{Global records and retained state}\label{sec:gobjects}
The global control state is
\[
 \mathcal S=(\gv,i,A,h,q,R,z,e_{\gv},p),
\]
where $A$ is the global view's activation proof, $R$ the adopted regional cursor,
$z$ the global no-data-progress count, $e_{\gv}$ its activation epoch, and $p$ either
empty or one pending rotation. A regional outcome always advances $h$; global
HOLD and CONFIRM need not do so. Initially these fields are fixed in $G_0$.
A control checkpoint at epoch $e$ is
\begin{equation}\label{eq:global}
 C_e=(c,e,g_{e-1},\mu,\mathcal S_e,\Hash(\mathbf r_e)),\qquad
 g_e=\Hash(\Tag{CHECKPOINT},C_e),
\end{equation}
where $\mu\in\{\Tag{ADVANCE},\Tag{HOLD},\Tag{ROTATE},\Tag{CONFIRM}\}$ and
$\mathbf r_e$ is the common report vector defined below. The predecessor is
$C_{e-1}$; even a no-progress epoch creates a new control record.

A report from $S\in V_j$ is
\[
 r_{S,e}=(c,e,g_{e-1},A_e,\mathcal R_{S,e},\mathsf{ready}),
\]
with an origin signature. $A_e$ is empty or the already existing finality proof
activating the parent's pending rotation. $\mathcal R_{S,e}$ is a prefix of the
regional outcome chain extending the effective state's cursor, with complete
new bodies and witnesses; at most $b_G$ outcomes and $B_G$ complete bytes are
admitted. A report may contain no extension. A true $\mathsf{ready}$ is the
signer's promise that it holds and retains the parent's entire adopted stream
and evidence and is prepared to join region $j$ at the next geographic
activation. It binds $(c,e,g_{e-1},\gv+1,j,h,q)$. Readiness is re-collected for each
parent; old acknowledgments cannot authorize a different tip.

Every non-bottom report must satisfy $\Proc{ValidReport}$: verify origin,
epoch and parent; validate the activation proof; validate the complete regional
extension under the effective global view; enforce size bounds and require its
bytes. Readiness alone never proves a range extending the parent. A signer may
lie about readiness, but a regional ready quorum contains honest data holders.
An honest node that accepts a complete report becomes a server for its body.
All honest members transmit true READY once its local condition holds; the
precise originating rule appears in \cref{alg:global}.

For external verification, $\Proc{ValidCheckpoint}(C_{e-1},C_e,\mathbf r_e)$
checks the epoch and parent hash, the report-vector commitment, every non-bottom
origin-signed report and its complete relative witness, and equality of the
action/state to \cref{alg:gtransition}'s deterministic result (including any
valid ancestor activation). Proofs may reference an already verified anchor;
a new receiver must obtain the missing ancestry. A report vector by itself is
not a portable agreement proof: a global vote certificate authenticates that
this checkpoint was endorsed after report agreement. The consumer checks both
the transition predicate and the FFG certificates; it need not reconstruct
private message-arrival times in the report protocol.

\subsection{Data transport and mathematical state}\label{sec:transport}
The logical data messages are
\begin{align*}
 &\Tag{ANNOUNCE}(c,\gv,\mathsf{objectID},h_0,h_1,\mathsf{kind}),\\
 &\Tag{GET}(c,\mathsf{objectID},[h_0,h_1]),\\
 &\Tag{DATA}(c,\mathsf{objectID},\mathsf{body},\mathsf{witness}).
\end{align*}
An object identifier binds a header, range and evidence type, not a peer address.
A supplier is free to be different from its original producer. A DATA response
is useful only if the entire requested object verifies against that identifier
and the receiver's certified context. Requests and replies need channel
authentication; their embedded producer statements and certificates remain
independently verifiable.


A sequencer persists its configuration, certified regional chain, accepted
global checkpoint chain, justified/finalized checkpoints, retained complete
objects and signing history. For the outstanding regional sequence it also
persists the highest retired retry, highest prepared certificate, and accepted
PREPREPARE and votes in each retry. Local timer origins are never reset to
re-enter a retired retry. A consumer persists its newest finalized control
anchor, highest activated global view, delivered log, preview suffix, authority
proofs and missing-object queue. These are mathematical states, not software
container requirements.

\section{The Protocol}\label{sec:protocol}

\subsection{Regional PBFT and transferable confirmation}\label{sec:regional}
The region runs a signed, weighted PBFT instance for each $I=(c,\gv,u)$.
PREPREPARE, PREPARE and COMMIT replace the previous broadcast/completion path.
Every member first authenticates a signed message and extracts any conflict
evidence before applying the phase-specific duplicate rules.

\begin{protocol}{alg:regional}{One regional PBFT sequence}
\textbf{Entry:} verify $R_{u-1}$ and its witness. For $u=1$, enter at the fixed
global activation boundary; otherwise enter when its predecessor certificate
first becomes available. Record entry time $a_S$ at member $S$ and start $v=0$.
A complete PREPREPARE may carry a missing predecessor certificate.
\begin{enumerate}[label=\textbf{R\arabic*.}]
\item \textbf{Propose.} In view zero, primary $\ell_i(u,0)$ authors a valid
batch with $\alpha=0$, signs PREPREPARE and multicasts it with its full witness.
In view $v>0$, use only the NEWVIEW procedure below: carry its selected batch,
or author at $\alpha=v$ if its selection is bottom. Push new mini-block bytes
to consumers; these bytes alone authorize no Preview.
\item \textbf{Prepare.} Accept the first valid PREPREPARE for the current
$(I,v,p_u)$ from its scheduled primary. Require the complete batch, $\alpha=0$
when $v=0$, and a valid matching NEWVIEW when $v>0$. Persist and multicast
PREPARE once for $(I,v)$. Never prepare a different digest in that retry.
The primary follows the same rule. Buffer an early proposal until local entry
into its sequence/retry; a future message alone never advances the retry.
\item \textbf{Commit vote.} On a matching PREPARE quorum for the accepted
proposal, persist $\PC_v(x)$ and the full batch, update the highest prepared
certificate, and multicast COMMIT once for $(I,v)$. Do so only while this
retry is active and has not been retired. COMMIT packets contain the signed
phase statement; they need not repeat the prepared certificate.
\item \textbf{Commit and export.} A complete valid $\RC_u$ for any retry,
including an older one, commits $R_u$. Persist it, relay the certificate once
to every regional member, and announce it to consumers. Relay only the first
accepted certificate for this semantic outcome; alternative signer encodings
do not trigger another relay. Bodies already cached need not be retransmitted.
A member can accept this proof without having voted in that retry.
\item \textbf{Advance.} Stop voting for sequence $u$ after committing. Start
$u+1$ only from this certified $R_u$, with retry zero and a new entry timer.
Continue serving older certificates and bodies. On learning a finalized
replacement global view, obey its activation and stop new old-view votes.
\end{enumerate}
\end{protocol}

\begin{protocol}{alg:localchange}{PBFT timeout, VIEWCHANGE, and NEWVIEW}
\textbf{Parameters:} $T_V\geq8\Delta_R$. Retry $v$ begins locally at
$a_S+vT_V$ while sequence $u$ remains unresolved.
\begin{enumerate}[label=\textbf{V\arabic*.}]
\item \textbf{Retire and report.} At entry to $v\geq1$, first persist the
retirement of all retries $<v$ and stop issuing their PREPARE/COMMIT messages.
Sign one immutable VIEWCHANGE to primary $\ell_i(u,v)$, reporting the highest
prepared certificate then known, or bottom. Include its complete witness.
Lower-view messages arriving later do not reopen voting or alter this signed
snapshot. A valid COMMIT certificate remains acceptable at any time.
\item \textbf{Construct.} At time $a_{\ell_i(u,v)}+vT_V+2\Delta_R$, an honest
primary that has not committed selects a fully available valid VIEWCHANGE
quorum. To make its choice deterministic, process available reports in sender
order and take the shortest prefix whose weight exceeds $2W_i/3$.
If no such quorum is available, do not propose; the next timeout still occurs.
Compute $\Proc{Select}(\NV_v)$, obtain its full batch, or build a fresh batch
if it is bottom, and multicast the signed NEWVIEW and PREPREPARE with all
report evidence and batch bodies. A primary signs at most one NEWVIEW per $(I,v)$.
\item \textbf{Validate.} A receiver verifies every report, its parent and
prepared certificate, the NEWVIEW primary and selection, and the complete
selected batch. Once in retry $v$, it may accept this proposal and run R2--R4.
Its own earlier prepared value may differ if the checked NEWVIEW selects a
higher prepared view; it does not impose a separate veto. A prepared
certificate alone is not a commit.
\item \textbf{Continue.} If the primary equivocates, withholds a required
object or fails to gather a quorum, keep the signing restrictions and retry at
the next local boundary. Do not invent a GAP, advance $u$, clear prepared
state, or interpret a timeout as a certificate of producer misbehavior.
\end{enumerate}
\end{protocol}

A recovered node restores its signing history before joining. It may receive
certificates and serve data immediately. If it missed a timer origin or cannot
establish safe current retry timing, it abstains until a later sequence's
certified predecessor or a new global activation. No historical votes are
backfilled. Regional liveness assumes a responsive quorum through the sequence;
a whole-region outage is recovered by the independent global layer.

The one-outstanding-sequence design makes the certified predecessor its stable
checkpoint. Retrying retains the same parent and carries at most one prepared
value per report. Portable prepared/commit certificates need no recursive
NEWVIEW history: quorum signers have already checked those justifications.
The next sequence is gated by commitment, but has no mandatory $n_i$-round
completion interval. Producer signatures remain those of the original author,
even when a different primary completes the prepared batch.

\begin{example}[A batch and a failed primary]
With ten equal-weight members, seven signatures form a quorum. The primary
proposes mini-blocks 101--110 after a certificate for 100. Seven PREPARE votes
form a prepared certificate and seven COMMIT votes certify the batch. If that
primary fails after enough members prepare, the next primary's VIEWCHANGE
quorum must preserve any batch protected by a possible COMMIT quorum. It
re-proposes the selected batch with its original signatures; it cannot replace
it merely because it prefers a different ordering.
\end{example}

\Needspace{34\baselineskip}
\subsection{Push delivery and missing-data repair}\label{sec:pull}
\begin{protocol}{alg:pull}{Stream transport and consumer acceptance}
\textbf{State at a receiver:} current delivered tip; a separate preview suffix;
certified authority; missing objects. A request is bounded to one configured
range. Larger gaps are divided into consecutive requests.
\textbf{Parameters:} announcement period $T_A$; request deadlines $\theta_R$
within a region and $\theta_G$ over the backbone.
\begin{enumerate}[label=\textbf{D\arabic*.}]
\item \textbf{Push and announce.} The active leader pushes each complete
mini-block and subsequently its regional certificate to subscribed consumers.
Holders announce their latest certified tip every $T_A$ and on a new certificate.
Each producer pushes each mini-block body at most once per subscribed consumer.
An announcement is a retrieval hint, not a confirmation event.
\item \textbf{Detect a gap.} A later position, an announced certified tip,
a checkpoint referencing unavailable bytes, or an expected-service timer may
reveal missing data. Request the exact missing header/body/certificate range
from the leader or a listed holder. A timer alone does not fabricate an empty
mini-block or replace the producer.
\item \textbf{Pull.} Try regional members in configured order without repetition;
after an unsuccessful attempt of $\theta_R$, try the next. For globally adopted
objects, also try the global holders via the backbone, with $\theta_G$ per attempt.
Each reply is capped at the requested size. Verify the complete response;
discard mismatched or incomplete objects and continue. After exhausting a list,
retry it after $T_A$. An honest holder responds with the requested retained
bytes and original evidence. Cache verified mini-block bodies by their content
digest across all proposal, report and certificate wrappers. Request only
missing bodies and metadata, with at most one outstanding attempt per body
at a receiver; cancel retries once a matching body arrives by push or pull.
An already in-flight response may still consume bandwidth and is counted below.
\item \textbf{Preview.} Check $\Proc{ValidBatch}$ and the full contiguous
regional outcome chain with $\Proc{VerifyRC}$ on every $\RC$. Check that the
activation proof authorizes the receiver's current global view and its fixed
activation time has been reached. Append verified positions in order
to the preview suffix and emit $\Tag{Preview}$. Missing predecessors prevent
output. Stop extending an old global view after learning its finalized successor.
\item \textbf{Deliver.} On a verified finalized global checkpoint covering a
new range, retrieve all missing bodies. Verify its finality links, every new
regional certificate and mini-block, exact predecessor continuity and the global
authority transitions. Append only previously undelivered positions and emit
$\Tag{Deliver}$ for each. Reject any attempt to replace the delivered prefix.
\item \textbf{Handoff.} Require the control proof to extend the receiver's
newest verified finalized anchor. An ancestor/replayed proof does not change
authority. Record a newer finalized anchor and process every not-yet-applied
finalized activation in its verified ancestry in increasing global view order,
including during catch-up. Queue each until its fixed activation boundary.
At each activation, adopt that rotation's retained tip and global view, keep all
delivered data and evidence, and remove preview entries beyond that retained tip.
Do not skip an intermediate activation or its preview truncation; missing
ancestry prevents advancing authority. Accept new producer data only under
the new activation proof.
Old signed statements remain eligible as historical evidence.
\end{enumerate}
\end{protocol}

The peer serving a range need not be its producer. Verification binds the
original producer and certified range, so changing a transport path changes
neither the accepted content nor its authorization. Repairing a packet loss
is separate from certifying a service failure. During normal operation the
proposer receives one primary full-data push, not a full copy from every voter.

\begin{figure}[htbp]\centering
\input{figures/stream-pipeline-pbft.tex}
\caption{The regional PBFT data path and the global control path. A regional
commit certificate authorizes Preview; global finality authorizes Deliver.
A proposal may contain several mini-blocks, and dependent sequences wait for
commitment. The diagram does not assign consumer execution or publication to
the sequencing service.}
\label{fig:stream}
\end{figure}

\begin{example}[A missing middle mini-block]
The consumer has 100 and receives 101 and 103 but not 102. It asks a regional
holder for 102 and any missing range certificate. A response with a different
body fails the committed digest check; a correct response restores the chain
101--102--103. The receiver then outputs the eligible contiguous range once.
Receiving 103 first never permits skipping position 102.
\end{example}


\subsection{Local authority and geographic recovery}
Local primary selection is fixed by \eqref{eq:primary} and the checked
NEWVIEW rule. A newly authored batch uses its current primary's producer key;
a carried batch preserves its original author view and key. Thus a local
primary change neither rewrites committed bytes nor reassigns authorship of a
prepared batch. Authentication of a batch is distinct from permission to
commit it: only the PBFT certificate establishes the latter.

A timeout cannot unlock a committed sequence. If a region loses its responsive
quorum, local progress can stop without producing any outcome. The global
protocol can nevertheless activate another region, or restart the recovered
current region in a fresh global view. That activation fixes a new $R_0$ and
new signing domains without waiting for the old unfinished sequence. It retains
every globally adopted entry; unadopted previews have only regional stability.

\subsection{Backbone synchronization and global report agreement}\label{sec:global}
Global epoch $e$ occupies $[(e-1)T_G,eT_G)$. Its control participants are all of
$V$, independent of the active production region. Logical signed control
messages use the backbone; complete bodies are pushed on the primary path or
pulled and cached as in \cref{alg:pull}. An honest recipient of a validated
report serves its body thereafter. The synchronized protocol consumes a message
only when this complete object is available.

\begin{protocol}{alg:global}{A complete global epoch}
\textbf{Common initial state:} $C_{e-1}$, its justification proof, the adopted
stream and any older finality proofs. $C_0$ is justified and finalized.
\begin{enumerate}[label=\textbf{G\arabic*.}]
\item \textbf{At epoch start.} If the parent has a pending rotation already
finalized by an earlier link, verify that proof and apply its activation once.
Each member snapshots the earliest contiguous known extension from the resulting
cursor that fits the configured bounds, taking the longest such prefix and
never truncating a certified batch. An honest member sets READY true exactly
when it holds the complete verified parent prefix and evidence and is ready
to start the next global view from that prefix in its configured region;
otherwise it sets false and
continues pulling the missing state. A healthy eligible member with these data
must set true. It signs one report; READY binds the parent's exact tip and
the next global view after any pre-transition activation. No extension is required.
\item \textbf{Report agreement.} Run $N$ parallel instances of
\cref{alg:ab}, with domain $(\Tag{REPORT},c,e,g_{e-1})$, round length $\Delta_G$
and validity predicate $\Proc{ValidReport}$. After $N\Delta_G$, replace every
missing or equivocated origin by $\botv$ to obtain $\mathbf r_e$.
No global target is selected before this completion.
\item \textbf{Derive the target.} Apply \cref{alg:gtransition} to the common
parent, activation proof and report vector. Construct exactly $C_e$ in
\eqref{eq:global}. Retain its complete adopted prefix and transition evidence.
\item \textbf{FFG vote.} At elapsed time $N\Delta_G$, each eligible member
signs $(\Tag{FFG},c,e-1,g_{e-1},e,g_e)$ after the checks in
\cref{alg:finality}. Disseminate the vote through the backbone.
\item \textbf{Certificate.} By $(N+1)\Delta_G$, responsive honest members
have a quorum of votes. They form and disseminate the link certificate.
By $(N+2)\Delta_G$, process justification and finality. Push or announce the
new checkpoint and proof to consumers; missing data remains retrievable.
\item \textbf{Boundary.} Set $T_G\geq(N+3)\Delta_G$; the remaining margin
makes the certificates and newly activated regional initialization available
before the next common boundary. Carry the adopted state to epoch $e+1$.
\end{enumerate}
\end{protocol}

Report agreement is an explicit, deliberately conservative convergence mechanism.
It already selects a unique control chain under synchrony. FFG supplies the
chosen two-link finality and activation rule on that chain; it is not invoked
as an unexplained solution to inconsistent reports. This construction pays for
report agreement even during a good epoch. It neither implements every Ethereum
fork-choice rule nor asserts Ethereum's performance for this different protocol.

A planned global view change becomes due $N_G$ checkpoint epochs after
activation; $K_G$ consecutive no-progress epochs trigger recovery.

\begin{protocol}{alg:gtransition}{Global state transition and region selection}
\textbf{Input:} common parent state, any finality-based activation applied at
epoch start, and $\mathbf r_e$.
\begin{enumerate}[label=\textbf{T\arabic*.}]
\item \textbf{Pending, not activated.} If a rotation remains pending, output
$\Tag{CONFIRM}$ with exactly the same stream tip, region and pending package.
No new regional data is adopted.
\item \textbf{Available extension.} Otherwise, among valid reported regional
chains rooted at the effective cursor, choose the longest one; break an exact
equality by its semantic digest. Let its endpoint be $(R^+,h^+,q^+)$, or retain
the parent cursor if no extension exists. These chains are compatible by
\cref{lem:rcunique}. Put $z^+=0$ if $h^+>h$, and $z^+=z+1$ otherwise.
\item \textbf{Eligible successor.} A region $j$ is ready if true READY
reports for the parent's exact tip have total weight $>2W_j/3$.
A change is due if $e-e_{\gv}\geq N_G$ or $z^+\geq K_G$.
If due and a different region is ready, choose the first such region after $i$
in the configured cyclic list. Otherwise, if $z^+\geq K_G$, no different region
is ready and $i$ is ready, choose $j=i$ for recovery. If a region is selected,
output $\Tag{ROTATE}$, retain the \emph{parent's} tip and cursor, and store
\[
 p=(\gv+1,j,h,q,R,e,\mathsf{readyProof},\mathsf{reason}).
\]
The successor starts with $u=1,v=0$ and its configured first primary. ROTATE does not
simultaneously adopt a new range: its readiness proof refers to the retained tip.
The proposed regional extension remains a preview. Selection of $j=i$ uses
the same finality and activation rules, with a fresh global view $\gv+1$ and $R_0$;
it does not reuse old instance identifiers or erase signing history.
\item \textbf{No rotation.} Otherwise adopt $(R^+,h^+,q^+)$, set the failure
count to $z^+$, and output ADVANCE if the data tip grows or HOLD otherwise.
If no region is selected, service can continue in the current global view.
\end{enumerate}
\end{protocol}

An honest reporter takes the earliest contiguous committed extension that fits
the bounds; it never drops an intervening outcome or truncates a certified batch.
Global no-progress means no new adopted mini-block, not low transaction volume
or dissatisfaction with a producer's ordering policy.

\subsection{Finality, activation and consumer verification}\label{sec:finality}
\begin{protocol}{alg:finality}{FFG links and geographic activation}
\textbf{Persistent state:} signed FFG history; justified and finalized checkpoints.
\begin{enumerate}[label=\textbf{F\arabic*.}]
\item \textbf{Before voting.} Verify the target derived by \cref{alg:global},
its complete adopted data and its parent justification. Sign at most one
statement for a target epoch. Never sign a pair with source/target epochs
$s_1<s_2<t_2<t_1$, nor a target conflicting with the known finalized prefix.
The concrete epoch protocol uses only $(e-1,e)$ links.
\item \textbf{Verify a link.} A link $C_s\Rightarrow C_t$ contains
$\QC_V(\Tag{FFG},c,s,g_s,t,g_t)$, the identified checkpoints and an ancestry
path. Require $s<t$, exact matching hashes and a justified source.
A valid link justifies its target.
\item \textbf{Finalize.} A justified $C_e$ is finalized by a link to its direct
child $C_{e+1}$. Ancestors are finalized too. An independently verifiable proof
contains the justification chain back to an already trusted anchor and this
adjacent link. Equal stream roots at different epochs do not mean equal
checkpoint identities.
\item \textbf{Activate.} If a pending ROTATE $C_r$ is finalized by the link
$C_r\Rightarrow C_{r+1}$, its successor starts at the fixed boundary beginning
epoch $r+2$, namely time $(r+1)T_G$. Before that boundary it verifies the finality
proof and retained data. The proof fixes $R_0$, global view, region and initial
leader for the new service. A node missing these objects does not start.
\item \textbf{Represent activation.} At the start of epoch $r+2$, apply this
proof to the strict ancestor ROTATE: set $(\gv,i)$ to its named successor,
$A=\pi_r$, $R=R_0$, $z=0$, $e_{\gv}=r+2$, and $p=\varnothing$;
preserve $(h,q)$. The first subsequent global target records this pre-transition
before evaluating \cref{alg:gtransition}. Checkpoints $C_r$ and
$C_{r+1}$ remain unchanged. New regional data cites the activation proof.
\end{enumerate}
\end{protocol}

At the activation boundary, every responsive honest global member has the
same proof and full retained data under the epoch budget. A healthy selected
region therefore starts from one common $R_0$, including after recovery in the
same region. An unavailable chosen region
cannot begin service; global epochs still run, detect the continued lack of
progress and attempt a later successor. A READY quorum proves possession at
its signing point, not immunity from a subsequent regional outage.

An old region disconnected from the control plane can continue to show signed
previews to stale clients. It cannot extend the accepted global history after
its replacement. All newly adopted data must match the effective global view
of the common global state. Consumers learn replacement from a finality proof,
not from an unauthenticated redirect or a private timeout.

\begin{example}[Local failure and geographic failure]
Suppose regional certificate $\RC_u$ covers position 110. If its primary stops,
VIEWCHANGE/NEWVIEW retries the next unfinished sequence from 110, preserving
any prepared batch required by its view-change proof. Now suppose the whole region becomes unreachable while the global
adopted tip is 100 and 101--110 exists only as regional preview. After $K_G$
no-progress epochs, a ready region can be named by ROTATE at $C_r$. CONFIRM at
$C_{r+1}$ finalizes that record, and the new region starts at epoch $r+2$ from
100. The consumer keeps every delivered position and drops only the unretained
preview. Had the global chain already adopted 110, the handoff would have to
retain 110, even if its checkpoint had not yet finalized when the outage began.
\end{example}

\Needspace{20\baselineskip}

\subsection{Equivocation and invalid inputs}\label{sec:evidence}
\begin{protocol}{alg:evidence}{Transferable conflict evidence}
\textbf{Input:} two origin-signed statements by the same configured member.
Verify both signatures, configuration, type and the comparison key below.
Return evidence exactly when the semantic contents differ under one key:
\begin{itemize}
\item ORDER: $(c,\gv,u,\alpha,h)$; BATCH: $(c,\gv,u,\alpha)$.
\item PREPREPARE, PREPARE or COMMIT: $(I,v)$ within that message type.
\item VIEWCHANGE or NEWVIEW: $(I,v)$ within that message type and sender.
Alternative witnesses for the same semantic report are not a conflict.
\item REPORT: $(c,e,g_{e-1},S)$ for the report's origin $S$.
\item FFG: equal target epochs and different source/target statements, or
source/target epoch pairs satisfying $s_1<s_2<t_2<t_1$ in either order.
\end{itemize}
Store and propagate at most one evidence pair per signer, type and comparison
key. Evidence never changes committee weights or thresholds within $c$.
\end{protocol}
Cross-retry votes are not automatically equivocation. Different certificates
for the same committed batch, later authorized authorship, relaying conflicting
origin values during global broadcast, and silence are not double-signing
proofs. Timeout only triggers recovery. Penalty amounts and adjudication lie
outside this service.

\section{Correctness and Soundness}\label{sec:proofs}
All results concern the fixed configuration and fault/delivery envelope in
Section~\ref{sec:model}. Cryptographic failures contribute a negligible error;
the remaining arguments are deterministic.

\begin{lemma}[Quorum intersection]\label[lemma]{lem:intersection}
Two quorums of a committee of weight $W_Q$ intersect in weight greater than
$W_Q/3$. Under a Byzantine bound below $W_Q/3$, they share an honest signer.
\end{lemma}
\begin{proof}
For signer sets $U,U'$, $\wt(U\cap U')\geq\wt(U)+\wt(U')-W_Q>W_Q/3$.
For an intersection with responsive honest members, use the bound on the union
of Byzantine and potentially unavailable members.
\end{proof}

\begin{lemma}[Authenticated origin broadcast]\label[lemma]{lem:ab}
In \cref{alg:ab}, all responsive honest members return the same value or bottom.
If the origin is honest and responsive, they return its valid input. Every
non-bottom output has a valid witness. The protocol terminates in $n$ rounds,
where $n$ is the committee size, if at least one member is honest and responsive
and complete-object delivery respects the round bound.
\end{lemma}
\begin{proof}
Consider a value retained by an honest responsive member in round $r$. If its
chain contains an honest responsive signer at position $k$, that signer sent
the length-$k$ prefix to all responsive members in round $k$. Every such member
receives it by that deadline and either retains it or already has two values.
If the chain has no honest responsive signer, its length is at most $n-1$.
The retaining member extends and forwards it for round $r+1$, unless it already
forwarded that value, in which case its earlier send suffices. A length-$n$
chain necessarily contains an honest responsive signer. Thus every value
retained honestly is, by round $n$, retained by every honest responsive member
or meets a receiver that already has two values.

If any honest responsive member has two values, apply this fact to both:
all such members end with two values and return bottom. Otherwise a nonempty
honest set propagates its sole value to all others, so all return it. If all
sets are empty, all return bottom. An honest origin signs only its own input;
unforgeability and first-round delivery exclude a second value or empty set.
Every retained value passed the fixed validity predicate, and fixed deadlines
ensure termination. The proof needs agreement on singleton-versus-multiple
outputs; cap-two sets need not contain the same pair when an origin sends three
or more different values.
\end{proof}


\begin{lemma}[Prepared uniqueness and carryover]\label[lemma]{lem:prepared}
For fixed $(I,v,p_u)$, two valid prepared certificates cannot name different
batches. If a batch $x$ has a valid COMMIT quorum in retry $v$, every valid
prepared certificate for $I$ in a later retry names $x$.
\end{lemma}
\begin{proof}
Two PREPARE quorums intersect in an honest member, which prepares only one
digest for $(I,v)$. For carryover, let $H_x$ be the honest signers of the COMMIT
quorum for $x$. Their weight exceeds $W_i/3$, and each persisted $\PC_v(x)$
before its COMMIT. Consider the earliest later retry $v'>v$ with a prepared
certificate for a different batch. At least one honest prepare signer checked
its NEWVIEW quorum. That VIEWCHANGE quorum intersects $H_x$ in an honest
member $S$. Member $S$ could not have sent its VIEWCHANGE for $v'$ before
issuing the older COMMIT: retirement would then prohibit that COMMIT.
It therefore reported a prepared certificate from a retry at least $v$.
By minimality of $v'$, all such certificates below $v'$ name $x$. The highest
reported prepared retry also lies in this interval and names $x$. The checked
NEWVIEW must select $x$, contradicting the honest PREPARE for a different
batch. The argument concerns signatures whenever assembled, so also covers
withheld COMMIT quorums revealed after a view change.
\end{proof}

\begin{theorem}[Regional PBFT safety]\label[theorem]{thm:bb}
Within a global view, all valid commit certificates for one sequence have the
same semantic outcome. Certified outcomes form one chain from authorized $R_0$,
and every committed batch has valid content and original producer authority.
These claims do not require a regional communication bound.
\end{theorem}
\begin{proof}
Within one retry, two conflicting COMMIT quorums require honest signers of
conflicting prepared certificates, excluded by \cref{lem:prepared}. Between
retries, the carryover part of that lemma excludes a conflicting prepared
certificate and hence any conflicting COMMIT quorum. Every COMMIT quorum
contains an honest signer that validated the full batch and its proposal
justification. The semantic $R_u$ depends on the batch and predecessor, not
on the retry that collected a certificate. Induction from $R_0$ proves a unique
outcome chain: honest members enter $u$ only after verifying its predecessor,
and all certificates for $u-1$ give the same predecessor. Hash binding and
producer signatures establish the range, ordering and author. Persistence
prevents restart from invalidating these signing restrictions.
\end{proof}

\begin{lemma}[Unique exported regional chain]\label[lemma]{lem:rcunique}
Valid regional extensions of the same activated $R_0$ are prefix-compatible,
including across local primary changes or after an interrupted retry.
\end{lemma}
\begin{proof}
Apply \cref{thm:bb} at the first sequence where two alleged chains differ.
A geographic activation has a new global-view identifier and therefore a new
$R_0$; the global protocol separately determines its retained prefix.
\end{proof}

\begin{theorem}[Regional PBFT progress]\label[theorem]{thm:regional-live}
Suppose the sequence starts from a common certified predecessor and a fixed
honest regional quorum remains responsive through entry and retries. If
$T_V\geq8\Delta_R$, a valid batch commits within
\begin{equation}\label{eq:slotbound}
 D_i=n_iT_V+7\Delta_R
\end{equation}
of the first responsive honest entry. In view zero, an honest primary's
complete proposal commits within $3\delta_R$ of its send.
\end{theorem}
\begin{proof}
Let $a=\min_{S\in H_i}a_S$. At activation all entries coincide. Otherwise the
first honest holder relays the predecessor certificate immediately, so all
responsive entries lie in $[a,a+\Delta_R]$. A proposal carrying that certificate
also permits entry; message order cannot prevent the recipient from joining.
The honest primary's view-zero proposal reaches all by one actual delay,
PREPARE votes by two, and COMMIT votes by three. Only commitment starts the
next sequence.

For retry $v\geq1$, all unresolved responsive members enter in
$[a+vT_V,a+vT_V+\Delta_R]$. If an honest member already committed and therefore
omits a VIEWCHANGE, its certificate reaches the others within $\Delta_R$ and
the sequence is done. Otherwise all honest VIEWCHANGE reports reach an honest
primary by $a+vT_V+2\Delta_R$. Its own scheduled construction time is no
earlier than this time and no later than $a+vT_V+3\Delta_R$. Complete honest
reports alone form a quorum. The primary selects from fully verified objects;
no missing Byzantine witness is required. Its NEWVIEW reaches all by
$a+vT_V+4\Delta_R$, PREPARE votes by $+5\Delta_R$, COMMIT votes by
$+6\Delta_R$, and the relayed certificate by $+7\Delta_R$. These events precede
the earliest next timeout. Early NEWVIEW messages are buffered, not rejected.
The cyclic schedule reaches a responsive honest primary within at most $n_i$
retries; \eqref{eq:slotbound} conservatively includes their preceding timers.
A selected old batch remains valid and contains at least one mini-block;
otherwise the honest primary supplies a new batch, allowing empty transaction
lists. No particular client's inclusion is needed.
\end{proof}

\begin{theorem}[Global agreement and authority]\label[theorem]{thm:global}
Under the global envelope, all responsive honest members derive the same $C_e$
in every epoch, justify it and finalize $C_{e-1}$. Adopted data prefixes are
monotone. A finalized ROTATE authorizes its unique named successor at the
specified boundary and preserves all previously globally adopted data.
\end{theorem}
\begin{proof}
Induct from genesis. The previous certificate phase makes the parent,
justification and required activation proof common before the epoch begins.
The report-validity predicate is consequently common. \Cref{lem:ab} supplies
the same report vector. The transition is deterministic; by
\cref{lem:rcunique}, its eligible regional extensions are compatible. Every
responsive honest member therefore derives the same target and retains its
complete adopted data.

Their weight exceeds $2W/3$, so their FFG votes establish the parent-to-child
link in one $\Delta_G$, and certificate dissemination finishes in another.
Both source and target epochs increase, satisfying honest signing constraints.
The target is justified and its parent finalized. ADVANCE extends the data
prefix; HOLD and CONFIRM retain it; ROTATE retains its parent's prefix and stores
one pending successor. Activation consumes the finality proof of that strict
ancestor and cannot mutate an existing checkpoint. This also holds when the
region is unchanged: the new global view separates all regional signing domains,
and its $R_0$ inherits the same retained global tip. Induction gives monotone
prefixes and a unique authorization chain. A conflicting certified target at
one epoch would also require an honest FFG double vote by quorum intersection.
\end{proof}
The FFG vocabulary follows Casper~\cite{ffg}. This proof establishes it on our
explicit synchronous control chain; it does not assume that FFG alone makes
inconsistent reports converge.

\begin{theorem}[Verification soundness and stream safety]\label[theorem]{thm:sound}
\Cref{alg:pull} satisfies the validity, integrity, agreement and verification
clauses of \cref{def:service}. Regional previews are prefix-compatible across
local changes. A global view change cannot remove a delivered entry.
\end{theorem}
\begin{proof}
An ORDER signature binds the producer to the header, and hash binding fixes
its ordered body and predecessor. \Cref{lem:rcunique} supplies one certified
regional chain with checked author and PBFT primary transitions.
\Cref{thm:global} supplies one global chain with monotone adopted prefixes and
unique geographic transitions. Finalized checkpoints lie on that chain.
Consumers verify these objects and append only complete, contiguous, previously
undelivered positions. Thus their delivered logs are prefixes of one sequence
and grow monotonically.

Changing contents or their order breaks a commitment; changing the parent,
position, view or global view changes a signed header. A substituted producer
requires either a valid authorization for that position or a forged signature.
A substituted confirmed history requires conflicting regional or global
certificates, excluded above. Serving the same object through a different peer
supplies none of these missing proofs. Correct complete evidence passes the
specified checks. At handoff, the retained tip includes every prior globally
adopted entry, hence every delivered entry. Applying every intervening activation
in order truncates an old preview at its first discarded suffix, even if a later
global view reuses those positions. Only an unretained preview is removed.
\end{proof}

\begin{lemma}[Data retention and repair]\label[lemma]{lem:pull}
A globally justified checkpoint has responsive honest holders of its complete
adopted prefix. For a bounded requested object, if $a$ attempted suppliers fail
before a responsive honest holder, retrieval takes at most $(a+1)\theta$ after
the first request, where $\theta$ is the complete attempt deadline. Regional
preview repair has the analogous bound during a live regional interval.
\end{lemma}
\begin{proof}
Responsive honest global members have weight $>2W/3$, possess the common adopted
prefix and retain it before voting. A verified vote quorum also intersects this
responsive set. Thus its data cannot exist solely at unavailable holders.
Each unsuccessful attempt costs at most $\theta$; an honest holder returns the
complete original object within that deadline. Its verification is independent
of the supplier being the producer. The local claim uses the responsive holders
after regional dissemination. Multiple chunks incur the sum of their bounds.
\end{proof}

\begin{theorem}[Service liveness]\label[theorem]{thm:live}
Assume the global envelope, a healthy eligible region that remains available,
and the capacity and minimum-tenure conditions in
\cref{sec:latency}. A finite sequence of failed local leaders is bypassed.
Following a regional outage and finitely many unproductive successor regions,
a healthy successor, possibly the recovered current region, resumes confirmation
and global delivery. During continuing
healthy service, live consumers' delivered logs grow without bound.
\end{theorem}
\begin{proof}
By \cref{thm:regional-live}, a responsive honest primary is reached and the
outstanding sequence commits within $D_i$. Retrying cannot discard a batch
protected by a COMMIT quorum. Each committed batch advances the stream by at
least one mini-block, even if all its transaction lists are empty.

If global progress stops, the common counter reaches $K_G$. Responsive members
of a healthy eligible region retain the global prefix and originate true READY
reports, even if an old PBFT sequence has no commit certificate. Their
weight gives a ready quorum. Report agreement delivers these reports without
completion of the failed instance. Distinct-region selections advance along
the cyclic list and cannot skip a continuously ready healthy region. If no
different region is ready, a ready current region is selected for recovery.
Thus after finitely many unsuccessful selections a healthy region is activated.
ROTATE and its confirming child finalize its authorization, and the clock fixes
the common start from a fresh $R_0$. The failed global view's decisions cannot
block these new instances. The minimum-tenure
condition allows a healthy successor to finish local recovery and report a new
batch before geographic guards replace it again.

Pull and report propagation supply that batch's data. The configured report
budget admits contiguous extensions and drains a bounded backlog. Subsequent
FFG links finalize adopted data; \cref{lem:pull} delivers it to live consumers.
Repeating this argument gives unbounded growth. It requires neither cooperation
from the failed region nor inclusion of a particular client transaction.
\end{proof}

\begin{proposition}[Sound equivocation evidence]\label[proposition]{prop:evidence}
Every accepted evidence pair in \cref{alg:evidence} identifies incompatible
statements signed by the named member. An honest member following the relevant
signing rules cannot produce such a pair.
\end{proposition}
\begin{proof}
Both signatures authenticate the exact comparison domain. Each listed predicate
violates its sign-once rule or the FFG surround rule. Different permitted retries,
views, configurations or relayed origin values do not by themselves satisfy
these predicates. Unforgeability excludes attributing an unsigned statement to
an honest member. This proves attribution of conflicting statements, not fault
attribution for an unanswered request or the size of a financial penalty.
\end{proof}

\section{Latency, Communication and Synchronization}\label{sec:complexity}
\subsection{Timing and capacity}\label{sec:latency}
For one complete bounded range, let
\begin{equation}\label{eq:pullbound}
 F_R=(a_R+1)\theta_R,\qquad F_G=(a_G+1)\theta_G,
\end{equation}
where $a_R,a_G$ count failed attempts before a usable holder. A useful announced
tip is discovered within $A_R=T_A+\Delta_R$ regionally and
$A_G=T_A+\Delta_G$ globally. For a consumer path use $A_P=T_A+\Delta_P$,
where $\Delta_P$ includes its announcement propagation. For $J$ sequential
chunks, use $\sum_{j=1}^J(a_j+1)\theta_j$. Let $S_G$ denote this bound for all
missing handoff state and $F_P$ the corresponding consumer repair bound.

Let $Q$ bound the report epochs needed to adopt a retrievable regional outcome
and all its still-unadopted predecessors while the global view stays active,
after an honest participating reporter holds those complete objects.
This is a workload bound, not a consequence of the maximum packing limits.
In a live epoch, an honest reporter includes the longest fitting contiguous
prefix, so at least one complete pending outcome can be adopted. If there are
at most $J_0$ pending outcomes up to a target when that reporter holds them,
$Q\leq J_0$ suffices: later outcomes are appended behind that target and cannot
overtake it. These are adoption epochs executing ADVANCE/HOLD, not epochs
spent on pending ROTATE/CONFIRM. Let $F_{\rm sync}$ be the sum of the per-chunk
global pull bounds for obtaining those objects; it equals $F_G$ for one bounded
request and is larger when several sequential requests are required.
A configured arrival envelope and adequate sustained byte service
can give a uniform $J_0$; otherwise use the actual backlog for that target.
The limits $b_G,B_G$ must permit each individual outcome and its relative
evidence in full. Nominal rates $b_G/T_G,B_G/T_G$ are only packing upper bounds;
they do not establish a service rate when outcomes have variable sizes.
With $D_i$ from \eqref{eq:slotbound}, define
\[
 L_{\rm sync}=A_G+F_{\rm sync}+QT_G,\qquad
 L_{\rm start}=\max_i D_i.
\]
To prevent repeatedly replacing a healthy newly activated region before it can
produce a reportable batch, require
\begin{equation}\label{eq:tenure}
 \min\{K_G,N_G\}T_G>L_{\rm start}+L_{\rm sync}+2T_G.
\end{equation}
The extra epochs cover report-cutoff alignment and the adopting epoch. This is
a conservative sufficient condition, not an optimal timer setting.

\subsection{Normal service and missing-data cases}

Let $L_{\rm build}$ bound deliberate batch accumulation and $L_{\rm queue}$
the wait for the previous sequence to commit. Once a view-zero primary sends
its complete proposal, the good-case bounds are
\begin{align}
 L_{\rm regional}&\leq3\delta_R,\label{eq:fast}\\
 L_{\rm preview}&\leq L_{\rm build}+L_{\rm queue}+4\delta_R+\delta_P.
 \label{eq:preview}
\end{align}
One delay carries PREPREPARE, one PREPARE and one COMMIT; a fourth conservatively
relays a full commit certificate. $\delta_P$ bounds the complete consumer path.
Batch construction must fit the admitted proposal interval for the regional
progress theorem; that theorem uses immediately eligible input, including empty
transaction lists. There is no forced $n_i$-round completion after success.
The reference does not pipeline dependent sequences: the next starts only on
a verified predecessor certificate. $3\delta_R$ is an upper latency bound, not
a lower bound from which a throughput ceiling can be inferred.

After regional certification, a conservative finalized-delivery bound in a
continuing healthy global view is
\begin{equation}\label{eq:deliver}
 L_{\rm final}\leq A_G+F_{\rm sync}+(Q+3)T_G+A_P+F_P.
\end{equation}
Besides backlog, this budgets one report-alignment epoch, one adopting epoch
and its confirming successor. An unadopted preview cancelled at a handoff has
no finalized-delivery promise. A committed batch carries at most $b_R$ mini-blocks.
Sustained throughput depends on accumulation, actual phase delays and capacity;
end-to-end delivery is also limited by global outcome and byte budgets.

\begin{table}[htbp]\centering\small
\begin{tabularx}{\textwidth}{@{}p{.33\textwidth}Y@{}}\toprule
Delivery case & Additional cost after data exists\\\midrule
Primary push succeeds & Primary data/certificate path; no consumer pull.\\
Known gap, first supplier honest & At most $\theta_R$ or $\theta_G$.\\
$a$ suppliers fail first & At most $(a+1)\theta$ for one bounded range.\\
Gap discovered by periodic announcement & Add the relevant $A_R,A_G$ or $A_P$.\\
$J$ missing chunks & Sum their attempt and verification bounds.\\
No retained reachable copy & No finite retrieval or lossless-recovery bound.\\\bottomrule
\end{tabularx}
\caption{Repair changes transport, not the confirmation predicate.}
\end{table}

\subsection{Local and geographic recovery}\label{sec:recoverybounds}

For $q$ unsuccessful PBFT retries followed by a responsive honest primary,
measured from the first honest entry of the unresolved sequence, a conservative
preview bound is
\begin{equation}\label{eq:localrecover}
 L_R(q)\leq qT_V+7\Delta_R+\delta_P.
\end{equation}
Add at most one $T_V$ when measuring from a fault inside an already running
retry, and $A_P+F_R$ for consumer discovery/repair. A retry with an honest
primary may commit a batch authored earlier; recovery here means stream
progress, not necessarily new authorship. A primary that succeeds intermittently
still commits data, while the base primary rotates after each committed batch.
For a newly activated region, $D_i$ bounds its first commitment from its
synchronized start. Local timers and prepared state are never cleared merely
to meet a latency estimate.

For a geographic outage or an interrupted regional instance requiring a fresh
global view, let $L_{\rm tail}$ bound the time until the last
previously produced but unadopted old-region range advances the global tip.
Such late data resets the no-progress counter, so this term cannot be silently
omitted. Let $T_{\rm ready}$ cover successor synchronization and READY-report
alignment; for a cold successor, $T_{\rm ready}\leq S_G+T_G$ once its inputs are
available. If the next chosen region is healthy, a conservative bound is
\begin{equation}\label{eq:georecover}
\begin{split}
 L_G\leq{}&L_{\rm tail}+T_G+K_GT_G+T_{\rm ready}+2T_G\\
          &{}+L_R(q)+A_P+F_P.
\end{split}
\end{equation}
The terms cover detection alignment, the no-progress guard, readiness,
ROTATE/CONFIRM and the fixed activation boundary, and first successor preview.
Readiness and guard formation may overlap; the bound does not subtract that
overlap. To obtain a new finalized Deliver, add \eqref{eq:deliver} for the
successor range. For a planned geographic handoff measured from its due boundary,
omit $L_{\rm tail}+T_G+K_GT_G$, retaining any readiness wait and the finality
stages. Each unsuccessful selected region adds another such detection and
transition interval. No finite bound follows without a global quorum, usable
honest backbone paths, retained bytes, transfer capacity and an eventual healthy
eligible successor. Recovery in the same region uses this bound and the same
READY reports and ROTATE/CONFIRM epochs. It adds no message type, quorum or
round to normal service; it still pays the global recovery delay.


\subsection{Control communication and verification}\label{sec:controlcost}
Let $n=n_i$ in a regional instance. Let $\kappa$ bound signature-plus-signer
bits, $L_R$ bound regional semantic metadata and relative predecessor evidence
excluding signatures and new bodies, and $L_G$ bound one global report's
semantic metadata and relative evidence excluding its newly appended chain
signatures and bulk bodies. These bounds include any missing proof ancestry;
checking a digest does not reconstruct it. A direct multicast counts as $n$
point-to-point transmissions. Let $E$ be the number of directed edges in the
connected backbone dissemination graph; a flooded logical message traverses
each admitted edge at most once. Body transfers use their actual route lengths.

PREPREPARE, all-to-all PREPARE and all-to-all COMMIT use $O(n^2)$ messages and
$O(n^2(L_R+\kappa))$ control bits in a stable view. Each commit certificate
contains $O(n)$ signatures. The conservative R4 rule lets every member relay
one full certificate to all peers, adding $O(n^2)$ messages and
$O(n^2L_R+n^3\kappa)$ bits. This cost is included even in the good case.
No aggregation or constant-size certificate is silently assumed.

Each VIEWCHANGE sent to the next primary contains at most one prepared
certificate: $O(L_R+n\kappa)$ control bits. NEWVIEW contains at most $n$ such
reports, so its multicast costs $O(n^2L_R+n^3\kappa)$ bits. Report delivery and
NEWVIEW use $O(n)$ logical point-to-point messages, in addition to the
$O(n^2)$ phase traffic. Our complete-object variant includes every reported
batch witness; without deduplication these retries add $O(n^2B_R)$ body bytes.
An honest primary admits at most one report per origin and chooses a quorum
of complete reports; an arbitrary stream of invalid packets is not admitted
protocol work. Each receiver processes at most one valid proposal/NEWVIEW per
retry for voting and one conflict pair per signing key. Capacity must cover
these bounds, including verification of all certificates in a NEWVIEW.

The unchanged global agreement runs $N$ signed broadcasts in parallel. In a
committee of size $k$, each origin's invocation forwards at most two semantic
values per member. Across $k$ origins there are $O(k^2)$ logical multicasts;
each signature chain has at most $k$ signatures. Hence
\begin{align}
 C_{\rm direct}(k,L)&=O(k^3(L+k\kappa)),\label{eq:directcost}\\
 C_{\rm backbone}(k,L)&=O(Ek^2(L+k\kappa)).\label{eq:backbonecost}
\end{align}
These are control-bit costs, with $O(k^3)$ direct or $O(Ek^2)$ backbone-edge
transmissions. Signature generation is $O(k^2)$. If $v_F$ bounds the verification
work of one detached validity witness, uncached endpoint checking is at most
$O(k^3(k+v_F))$ operations; backbone forwarding can add
$O(Ek^2(k+v_F))$. Previously checked witnesses can be cached by semantic digest.

\begin{table}[htbp]\centering\small
\begin{tabularx}{\textwidth}{@{}p{.30\textwidth}p{.18\textwidth}Y@{}}\toprule
Component & Transmissions & Control bits\\\midrule
Regional PBFT phases & $O(n_i^2)$ & $O(n_i^2(L_R+\kappa))$\\
Regional certificate relay & $O(n_i^2)$ & $O(n_i^2L_R+n_i^3\kappa)$\\
One VIEWCHANGE/NEWVIEW & $O(n_i)$ & $O(n_i^2L_R+n_i^3\kappa)$\\
Global report agreement & $O(EN^2)$ & $O(EN^2(L_G+N\kappa))$\\
Global FFG and link export & $O(EN)$ & $O(ENL_G+EN^2\kappa)$\\\bottomrule
\end{tabularx}
\caption{Unaggregated reference costs. Regional rows count direct transmissions;
global rows count backbone-edge transmissions. Bulk bodies are additional.}
\end{table}

With $q$ unsuccessful retries followed by success, regional phase, view-change
and export traffic is $O((q+1)n_i^2)$ messages and
$O((q+1)(n_i^2L_R+n_i^3\kappa))$ control bits. Under a continuously live region,
the cyclic schedule reaches an honest responsive primary after at most
$n_i-1$ preceding identities. Weight below one third does not imply that fewer
than one third of the identities are faulty or unavailable. Consequently the
worst first-success bound can still be cubic in $n_i$ in total messages; the
improvement is quadratic normal-view traffic and the removal of compulsory
cubic authenticated-broadcast completion from \emph{every} successful batch.
The old completion's $O(n_i^3(L_R+n_i\kappa))$ control-bit term is absent.
The global signed-broadcast term is unchanged and may dominate system cost.

A normal phase receiver checks $O(n_i)$ signatures; verifying one full
NEWVIEW or all uncached certificate variants can cost $O(n_i^2)$ signatures.
This is separate from batch execution, which this service does not perform.
Over duration $T$ with $U$ new-certificate events, the all-holder announcement
policy costs $O(N\lceil T/T_A\rceil+NU)$ logical messages; multiply flooded
backbone announcements by at most $E$. For $J$ bounded pull requests with at
most $a$ failed suppliers each, requests/replies add $O((a+1)J)$ messages and
their path lengths. $Z$ forwarded conflict proofs add $O(EZ)$ transmissions
and their actual proof sizes. These interval costs and consumer subscriptions
are separate from per-batch agreement.

\subsection{Payload and retained state}\label{sec:payloadcost}
Let $X$ be the distinct body bytes requested in an interval, including speculative
objects later discarded, and $M$ the sum of missing bytes over $R$ receivers.
Then $M\leq RX$; if one of the $R$ receivers already holds everything,
$M=(R-1)X$ when all others need the full range. A transfer plan sending each
needed content digest once per receiver uses $M$ body bytes. The primary path,
however, pushes per mini-block and may repeat identical bodies at different
positions. Define $M_0$ as the sum of all primary-push body sizes across their
destinations: normal loss-free service with no repair costs $M_0$ body bytes.
Let $M_r$ be the sum of still-missing
body sizes entering repair. With at most $a$ failed attempts per repaired body,
each response capped at its requested length,
\begin{equation}\label{eq:payload}
 B_{\rm body}\leq M_0+(a+1)M_r\leq(a+2)M.
\end{equation}
The last inequality assumes $M_0,M_r\leq M$, as when each receiver has at most
one primary push and one sequential repair campaign per initially missing body.
Pure pull with no primary push gives $(a+1)M$. A late primary push racing a
successful pull explains the extra copy; cancellation cannot retract in-flight
bytes. Body caching is shared across all origin-report wrappers, not just
across repetitions of the same report identifier.
Announcements, requests and certificates are additional control traffic.
Physical link-byte cost multiplies each transfer by its route length; it is
not the $E$ factor for a flooded compact control envelope.

The leader sends one primary full copy to $P$ while regional validators obtain
the proposal for voting. Global report verification can require every responsive
global member to fetch newly adopted bodies. Backbone/pull avoids copying full
bodies along every logical vote edge, but does not eliminate those required
replicas. If caching is disabled, conservative complete-object bounds are
$O(n_i^2(B_R+L_R+n_i\kappa))$ bits per regional retry, including
uncached certificate/body relay and full NEWVIEW bundles, and
$O(EN^2(B_G+L_G+N\kappa))$ for global report exchange. In a healthy view-zero
batch, shared body caching reduces regional body replication to $O(n_iB_R)$
bytes, while the unaggregated certificate-relay bit term remains cubic.

A full replica retains $O(8X_{\rm history}+A_{\rm history})$ bits, where $X$ is
body bytes and $A$ is evidence bits. Constant byte/bit conversions are suppressed
in the preceding asymptotic bounds. A cold handoff pays for
both missing data and missing evidence; a hot handoff may fetch only a new
authorization proof. No constant-size root reconstructs unavailable bytes.
Pruning requires a separately specified archive/retrieval rule; the reference
retains the configuration's history.


\subsection{Limits and security interpretation}\label{sec:limits}
Regional certificate safety survives arbitrary message delay under the fixed
fault bound and persistent-state assumptions. Its timer-based progress proof
requires the stated regional interval and complete-object capacity. The global
report protocol still requires continuous synchronous rounds: violating its
round bounds removes its agreement guarantee. In particular, restoring a
network after inconsistent global target votes does not by itself supply an
unspecified stalled-epoch recovery procedure. Regional PBFT does not repair
that global-layer limitation.

If responsive global weight is at most $2W/3$, finality and geographic recovery
may stop even when a successor region is locally healthy. A failed region may
therefore prevent its own replacement by removing too much global weight.
Prepared data withholding is resisted within a live region because honest
signers retain full bytes; this gives no retrieval guarantee after all such
holders become unreachable. Unbounded backlog or insufficient tenure also
invalidates the progress bounds. None of these failures can be hidden by
counting only message headers. Execution correctness, transaction fairness,
financial penalties and service under arbitrary partitions are outside the
proved properties.

\section*{Relation to Existing Protocols}
The regional construction instantiates the PREPREPARE/PREPARE/COMMIT and
prepared-certificate view-change structure of PBFT~\cite{pbft}. Weighted
quorums, a single outstanding batch, retry counters reset per sequence, fixed
bounded-skew retry timing and immutable original-author metadata are explicit
choices of this specification. The safety and progress arguments above prove
these choices rather than importing the original implementation's performance
or its full partial-synchrony machinery. The global layer retains signed-chain
broadcast following Dolev and Strong~\cite{ds}, and adjacent FFG-style finality
links~\cite{ffg}. Authenticated global report agreement selects its target
before FFG voting. All necessary algorithms, validity predicates and proofs
are given here; the cited protocols identify the designs' origins.

\appendix
\section{Complete Authenticated Broadcast}\label{app:ab}
Fix a committee $Q$ of $n$ identities, origin $o\in Q$, common instance domain
$J$, deterministic complete-object validity predicate $F_J$, and round length
$\Delta$. The special input value $\botv$ has an empty witness. Invalid or absent
traffic is not a signed bottom input, although the algorithm's bottom output
can represent an absent, ambiguous or genuinely bottom-valued origin. The
calling protocols intentionally treat these outputs identically.

For semantic value $z$, a length-$r$ chain is
\[
 C_r=((p_1,\sigma_1),\ldots,(p_r,\sigma_r)),\quad p_1=o,
\]
\[
 \sigma_k=\Sig_{p_k}(\Tag{CHAIN},J,o,z,k,C_{k-1}),\qquad C_0=().
\]
$\Proc{ChainValid}$ checks $1\leq r\leq n$, distinct configured identities,
first signer $o$, and each nested signature. The detached witness $\omega_z$
must satisfy $F_J(z,\omega_z)$. Different valid witnesses are not different
semantic values. Receivers obtain the complete required body before accepting;
honest forwarding members become servers for it.

\begin{protocol}{alg:ab}{Signed synchronous broadcast of one origin's value}
\textbf{State at $p\in Q$:} retained set $A_p=\varnothing$, of at most two
semantic values; forwarded set $F_p=\varnothing$. Round $r$ has send boundary
$(r-1)\Delta$ and receive deadline $r\Delta$; there are $n$ rounds.
\begin{enumerate}[label=\textbf{A\arabic*.}]
\item \textbf{Origin at round 1.} For valid input $z$, send
$(J,o,z,C_1,\omega_z)$ to all members, including itself. Mark $z$ forwarded
at the origin. Self-delivery is immediate.
\item \textbf{Receive.} Verify the chain and complete witness. Associate a
length-$r$ chain only with round $r$; buffer early arrivals and ignore an object
not completely valid by $r\Delta$. Enforce configured limits and admission
quotas while reserving every prescribed honest message.
\item \textbf{At deadline $r\Delta$.} Process valid length-$r$ packets in
semantic-digest order. Retain $z$ and its witness if already retained or if
$|A_p|<2$; ignore other values for this call. If $r<n$, $z$ is retained but not
forwarded, and some received valid chain for $z$ omits $p$, append $p$'s signature
and multicast the length-$(r+1)$ chain at this boundary. Mark it forwarded.
Forward at most once per retained value per origin over the whole invocation.
Continue prescribed forwarding after retaining two values.
\item \textbf{At deadline $n\Delta$.} After processing, return the sole element
of $A_p$ if $|A_p|=1$, otherwise return $\botv$. Retain a non-bottom witness.
\end{enumerate}
\end{protocol}

A round-$r$ send reaches responsive honest receivers by $r\Delta$ by the
complete-object premise. If a chain already contains $p$, its earlier prefix
was signed and disseminated by $p$, so $p$ does not append itself twice.
Parallel calls have independent origin identifiers and cap-two state. For
global reports, $F_J$ accepts bottom or an origin-signed valid report
for the fixed parent and epoch. The round count is $n$, not a count inferred
from the stake threshold. \Cref{lem:ab} proves this concrete algorithm.


\begin{thebibliography}{9}\small
\setlength{\itemsep}{2pt}
\bibitem{pbft} M. Castro and B. Liskov.
\emph{Practical Byzantine Fault Tolerance}. OSDI 1999, pp. 173--186,
Sections 4.2 and 4.4.
\url{https://www.usenix.org/legacy/events/osdi99/full_papers/castro/castro_html/node4.html}.
\bibitem{ffg} V. Buterin and V. Griffith.
\emph{Casper the Friendly Finality Gadget}. 2017.
\url{https://arxiv.org/abs/1710.09437}.
\bibitem{ds} D. Dolev and H. R. Strong.
\emph{Authenticated Algorithms for Byzantine Agreement}.
SIAM Journal on Computing 12(4):656--666, 1983.
\url{https://doi.org/10.1137/0212045}.

\end{thebibliography}
\end{document}
